\scanpage{10}
\deleted{On peut supposer que \inserted{$k$ alg. clos} $X,Y$ affines.}
\margin{\unsure{!}}

Equivaut à ceci : \deleted{\unsure{de|$k$}} Si $u(x,\deleted{\unsure{a}}b)=1$, $u(\unsure{-|a},y)=1$ \unsure{identiqut}, alors $u(x,y)=1$ \unsure{identiqut}. Sous cette forme, on peut supposer \deleted{\unsure{à|$k$}} $k$ alg. clos.

\deleted{Alors \inserted{\unsure{Rien n'est|Bien entendu} \unsure{changé}} l'énoncé \unsure{précédent|équivalent} \unsure{consiste à|aussi à} \unsure{Notons que l'énoncé|N.B. que l'énoncé} \unsure{à remplacer|implique} $X,Y$ par Alors l'énoncé \unsure{vient de cet énoncé} : \unsure{équivaut|génériquement} à : \unsure{tout} \inserted{$\varphi(x,y)$} \unsure{hom} $X\times Y\unsure{\mathord{-}|\to}G_m$ \unsure{a la forme|est de la forme}}

On peut supposer $X,Y$ affines, et une réduction immédiate nous ramène au cas où $X,Y$ de dim 1.

\uline{Corollaire} Kif-kif, $G_m$ remplacé par un \deleted{\unsure{tore|des|deux}} \unsure{gpe|groupe} de \unsure{t. m.}

\deleted{$X\times Y\subset\overline X\times\overline Y$}
\deleted{$\unsure{\mathcal O|O}(\overline X\times\overline Y)$}

\uline{Corollaire} 1 \inserted{Soit $T$ un tore.} Tout morphisme \deleted{d'un} $u:G\mathrel{-}G_m$, \deleted{qui} $G$ \unsure{gpe|groupe} \uline{lisse} \uline{connexe}, qui envoie $e$ dans $e$, est un hom de groupes. \deleted{\unsure{De plus|D'ailleurs}}

\deleted{\unsure{Soit}} \deleted{\uline{Corollaire} 3} \deleted{\uline{Proposition} \unsure{2|3|?} $\mathrm{Hom}(G,G_m)$ est \unsure{un groupe de type fini.}} \deleted{\unsure{Prop}} $\mathrm{Hom}_k(G,G_m)$ \deleted{Déf} est un gpe de type fini.
% archive sheet number: 9
